\documentclass{jsarticle}
\usepackage{amsmath}
\usepackage{amssymb}
\begin{document}
\title{瞑想法と速読法との関係、そして \\
       見て分かるという状態をつくる}
\author{}
\date{}
\maketitle



    瞑想法は、側頭葉を緩めてリラックスさせると、頭頂葉も楽にリラックスできる。
  このリラックスした状態で何も考えない状態を維持すると瞑想法になる。
  この瞑想法で速読法に必要な視野の確保と拡大、そして潜在意識に楽に情報が
  入るので、情報を文字として分散入力が容易になり、意味もすんなり分かる。


    GSR速読法も頭頂葉から下半身にかけて、気の流れを生成させて瞑想法をすると、
  情報が文字として自然と分かる。リラックスする状態になると、潜在意識に
  情報を入力する事ができる。このgenarative state瞑想は気分が楽になる。


    瞑想法は側頭葉と頭頂葉をリラックスさせて、情報としての文字を潜在意識に
  入力する役割がある。

\end{document}
